%% toc: P15 164 164.7,164.8
% Proofs for the reconstructions accompanying the Eberhardt--Glymour history of Reichenbach's
% probability logic and Niiniluoto's history of the Hintikka program.
\begin{annexitem}{annex:p15-w02b-p15-030-flat-prior}{thm:p15-w02b-p15-030-flat-prior}{What a flat prior does and does not reproduce}
\begin{proof}
\emph{Beta integrals.} For integers \(a,b\geq0\) put
\(J(a,b)=\int_0^1\theta^a(1-\theta)^b\,d\theta\). We show
\(J(a,b)=a!\,b!/(a+b+1)!\) by induction on \(b\). For \(b=0\),
\(J(a,0)=1/(a+1)=a!\,0!/(a+1)!\). For \(b\geq1\), integration by parts with
\(u=(1-\theta)^b\), \(dv=\theta^a d\theta\) gives
\[
 J(a,b)=\Bigl[\frac{\theta^{a+1}}{a+1}(1-\theta)^b\Bigr]_0^1
 +\frac{b}{a+1}\int_0^1\theta^{a+1}(1-\theta)^{b-1}d\theta
 =\frac{b}{a+1}J(a+1,b-1),
\]
the boundary term vanishing at both ends. By the induction hypothesis
\(J(a+1,b-1)=(a+1)!\,(b-1)!/(a+b+1)!\), so
\(J(a,b)=a!\,b!/(a+b+1)!\).

\emph{Posterior.} Given \(\theta\), the probability of any particular
sequence of \(n\) outcomes with \(k\) successes is
\(\theta^k(1-\theta)^{n-k}\). With prior density \(1\) on \([0,1]\), Bayes's
formula for densities gives posterior density
\(\theta^k(1-\theta)^{n-k}/J(k,n-k)\), which is proportional to
\(\theta^k(1-\theta)^{n-k}\).

\emph{Mode.} Let \(g(\theta)=\theta^k(1-\theta)^{n-k}\). If \(k=0\), \(g\) is
strictly decreasing on \([0,1]\) (as \(n\geq1\)) and its unique maximizer is
\(0=k/n\). If \(k=n\), \(g=\theta^n\) is strictly increasing and the maximizer
is \(1=k/n\). If \(0<k<n\), \(g\) vanishes at \(0\) and \(1\) and is positive
inside; on \((0,1)\), \(\log g=k\log\theta+(n-k)\log(1-\theta)\) has derivative
\(k/\theta-(n-k)/(1-\theta)\), which is positive for \(\theta<k/n\), zero at
\(\theta=k/n\), and negative for \(\theta>k/n\). So \(k/n\) is the unique
maximizer.

\emph{Mean.} The posterior mean is
\(J(k+1,n-k)/J(k,n-k)=\frac{(k+1)!\,(n-k)!}{(n+2)!}\cdot\frac{(n+1)!}{k!\,(n-k)!}=\frac{k+1}{n+2}\).
Finally \((k+1)/(n+2)=k/n\iff n(k+1)=k(n+2)\iff n=2k\iff k/n=1/2\). With
\(n=k=1\) the mode is \(1\) and the mean \(2/3\).
\end{proof}
\end{annexitem}

\begin{annexitem}{annex:p15-w02b-p15-030-russell-primes}{thm:p15-w02b-p15-030-russell-primes}{Russell's rearrangement of the integers}
\begin{proof}
Let \(\Pi\) be the set of primes and \(\Pi^c\) the other positive integers;
both are infinite (there are infinitely many primes, and every even number
greater than \(2\) is in \(\Pi^c\)). In any enumeration, write \(P_N\) for the
number of primes among the first \(N\) terms.

\emph{Case \(0<r<1\).} Build \(m_1,m_2,\ldots\) as stated: at step \(N+1\)
(\(N\geq0\)), if \(P_N<rN\) take the least prime not yet used, otherwise the
least non-prime not yet used. We show \(|P_N-rN|\leq1\) for all \(N\). For
\(N=0\), \(P_0=0\). Suppose \(|P_N-rN|\leq1\). If \(P_N<rN\), then
\(P_N-rN\in[-1,0)\) and \(P_{N+1}-r(N+1)=(P_N-rN)+1-r\in[-r,1-r)\subseteq[-1,1]\).
Otherwise \(P_N-rN\in[0,1]\) and \(P_{N+1}-r(N+1)=(P_N-rN)-r\in[-r,1-r]
\subseteq[-1,1]\). Hence \(|P_N/N-r|\leq1/N\to0\).

Each class is drawn from infinitely often. If only finitely many primes
were taken, \(P_N\) would be bounded while \(rN\to\infty\), contradicting
\(P_N\geq rN-1\). If only finitely many non-primes were taken, say \(c\),
then \(P_N=N-c\) for large \(N\) and \(P_N-rN=(1-r)N-c\to\infty\),
contradicting \(P_N\leq rN+1\). Since each class is taken in increasing
order and infinitely often, every prime and every non-prime eventually
appears, each exactly once; so the sequence is an enumeration of the
positive integers.

\emph{Case \(r=0\).} Put the primes, in increasing order, at the square
positions \(1,4,9,\ldots\) and the non-primes, in increasing order, at the
other positions. This is an enumeration, and \(P_N=\lfloor\sqrt N\rfloor\), so
\(P_N/N\leq N^{-1/2}\to0\).

\emph{Case \(r=1\).} Exchange the roles: non-primes at square positions,
primes elsewhere. Then \(N-P_N=\lfloor\sqrt N\rfloor\) and \(P_N/N\to1\).
\end{proof}
\end{annexitem}

\begin{annexitem}{annex:p15-w02b-p15-030-ravens}{thm:p15-w02b-p15-030-ravens}{Probabilistic implication does not contrapose}
\begin{proof}
Since \(P(\neg B)>0\),
\(P(\neg R\mid\neg B)=1-P(R\land\neg B)/P(\neg B)\), and
\(P(R\land\neg B)=P(R)P(\neg B\mid R)=P(R)(1-P(B\mid R))\) because
\(P(R)>0\). This is the displayed identity.

For the example, \(P(R\land\neg B)=0.9\times0.1=0.09\) and
\(P(\neg R\land\neg B)=P(\neg R)(1-P(B\mid\neg R))=0.1\times0=0\), so
\(P(\neg B)=0.09\) and \(P(\neg R\mid\neg B)=1-0.09/0.09=0\). The hypotheses
\(0<P(R)<1\) and \(P(\neg B)>0\) hold. Replacing \(0.9\) by
\(1-\eta\) in both places, with \(P(B\mid\neg R)=1\), gives
\(P(B\mid R)=1-\eta\) arbitrarily close to \(1\) and still
\(P(\neg R\mid\neg B)=0\). The final sentence is read off the identity:
\(P(\neg R\mid\neg B)\geq1-\epsilon\) exactly when
\(P(R)(1-P(B\mid R))\leq\epsilon P(\neg B)\).
\end{proof}
\end{annexitem}

\begin{annexitem}{annex:p15-w02b-p15-030-lambda}{thm:p15-w02b-p15-030-lambda}{Elementary properties of the representative function}
\begin{proof}
(1) For \(n\geq1\),
\(\frac{n}{n+\lambda}\cdot\frac{n_i}{n}+\frac{\lambda}{n+\lambda}\cdot\frac1K
=\frac{n_i+\lambda/K}{n+\lambda}\). Hence
\[
 \frac{n_i+\lambda/K}{n+\lambda}-\frac{n_i}{n}
 =\frac{\lambda}{n+\lambda}\Bigl(\frac1K-\frac{n_i}{n}\Bigr),
\]
and since both \(1/K\) and \(n_i/n\) lie in \([0,1]\), the bracket has
absolute value at most \(1\). The bound \(\lambda/(n+\lambda)\) does not depend
on \(n_i\) and tends to \(0\) as \(n\to\infty\), which gives (3) along every
sequence of samples.

(2) The denominators are positive. Cross-multiplying,
\((n_i+1+\lambda/K)(n+\lambda)>(n_i+\lambda/K)(n+1+\lambda)\) is equivalent,
after cancelling \((n_i+\lambda/K)(n+\lambda)\), to
\(n+\lambda>n_i+\lambda/K\), that is \((n-n_i)+\lambda(1-1/K)>0\). This holds
because \(n_i\leq n\), \(\lambda>0\) and \(K\geq2\).

(3) For \(n\geq1\), as \(\lambda\to0\) numerator and denominator tend to
\(n_i\) and \(n\). As \(\lambda\to\infty\), dividing by \(\lambda\) gives
\((n_i/\lambda+1/K)/(n/\lambda+1)\to1/K\). Putting \(\lambda=K\) gives
\((n_i+1)/(n+K)\), which is (5); with \(n_i=n\) and \(K=2\) it is
\((n+1)/(n+2)\), which is (6).
\end{proof}
\end{annexitem}

\begin{annexitem}{annex:p15-w02b-p15-030-principal}{thm:p15-w02b-p15-030-principal}{Von Wright's Principal Theorem of Confirmation (7)}
\begin{proof}
Since \(h\) entails \(e_n\&i_{n+1}\), it entails \(e_n\), so
\(P(h\&e_n)=P(h)\) and \(P(h\&e_n\&i_{n+1})=P(h)\). Also
\(P(e_n)\geq P(e_n\&i_{n+1})>0\). Then
\[
 P(h/e_n\&i_{n+1})=\frac{P(h)}{P(e_n\&i_{n+1})}
 =\frac{P(h)/P(e_n)}{P(e_n\&i_{n+1})/P(e_n)}
 =\frac{P(h/e_n)}{P(i_{n+1}/e_n)}.
\]
Here \(P(h/e_n)=P(h)/P(e_n)>0\) because \(P(h)>0\), and
\(0<P(i_{n+1}/e_n)<1\), so the quotient exceeds \(P(h/e_n)\).
\end{proof}
\end{annexitem}

\begin{annexitem}{annex:p15-w02b-p15-030-likelihood}{thm:p15-w02b-p15-030-likelihood}{The likelihoods (16) and (16\('\)) from the predictive rule}
\begin{proof}
Let \(e\) state, for \(m=1,\ldots,n\), that \(a_m\) is in \(Q_{k(m)}\), each
\(Q_{k(m)}\) among the kinds of \(C^w\). Write \(e_m\) for the description of
\(a_1,\ldots,a_m\) and \(m_j\) for the count of \(Q_j\) in \(e_m\). By the
product rule and the hypothesis,
\[
 P(e/C^w)=\prod_{m=0}^{n-1}P(Q_{k(m+1)}(a_{m+1})/e_m\&C^w)
 =\prod_{m=0}^{n-1}\frac{m_{k(m+1)}+\lambda/w}{m+\lambda}.
\]
The denominators multiply to \(\prod_{m=0}^{n-1}(m+\lambda)=\Gamma(n+\lambda)/\Gamma(\lambda)\),
by repeated use of \(\Gamma(x+1)=x\Gamma(x)\). In the numerators, the factors
contributed by the successive members of a fixed observed kind \(Q_j\) are
\(\lambda/w,\,1+\lambda/w,\ldots,(n_j-1)+\lambda/w\), whose product is
\(\Gamma(n_j+\lambda/w)/\Gamma(\lambda/w)\); kinds not observed contribute no
factor. Multiplying over the \(c\) observed kinds gives (16).

With \(\lambda=w\): \(\Gamma(w)/\Gamma(n+w)=(w-1)!/(n+w-1)!\) and
\(\Gamma(n_j+1)/\Gamma(1)=n_j!\), which is (16\('\)).
\end{proof}
\end{annexitem}

\begin{annexitem}{annex:p15-w02b-p15-030-combined}{thm:p15-w02b-p15-030-combined}{Posterior behaviour of the generalized combined system}
\begin{proof}
Write \(G(w)=\Gamma(\alpha+w)/\Gamma(w)\) for the unnormalized prior weight
of a constituent of width \(w\geq1\), and \(Z\) for the normalizer, so
\(P(C^w)=G(w)/Z\). A constituent incompatible with \(e\) has likelihood \(0\).
For a compatible \(C^w\), (16\('\)) gives
\(P(C^w)P(e/C^w)=Z^{-1}\,\Gamma(\alpha+w)\,\Gamma(n+w)^{-1}\prod_j n_j!\).

(1) The compatible constituents are those with \(CT\) containing the \(c\)
observed kinds: \(\binom{K-c}{i}\) of width \(c+i\) for \(i=0,\ldots,K-c\). By
Bayes's formula (12) the common factor \(Z^{-1}\prod_jn_j!\) cancels, and
\[
 P(C^w/e)=\frac{\Gamma(\alpha+w)/\Gamma(n+w)}
 {\sum_{i=0}^{K-c}\binom{K-c}{i}\Gamma(\alpha+c+i)/\Gamma(n+c+i)},
\]
which is (20) with \((x-1)!=\Gamma(x)\). If \(n=\alpha\), each ratio
\(\Gamma(\alpha+x)/\Gamma(n+x)\) equals \(1\), and the value is
\(1/\sum_i\binom{K-c}i=1/2^{K-c}\).

(2) For \(i\geq1\),
\[
 \frac{\Gamma(\alpha+c+i)/\Gamma(n+c+i)}{\Gamma(\alpha+c)/\Gamma(n+c)}
 =\prod_{t=0}^{i-1}\frac{\alpha+c+t}{n+c+t}
 \leq\Bigl(\frac{\alpha+K}{n+c}\Bigr)^{i},
\]
which tends to \(0\) as \(n\to\infty\). Dividing numerator and denominator of
(20) at \(w=c\) by its own numerator,
\[
 P(C^c/e)=\Bigl(1+\sum_{i=1}^{K-c}\binom{K-c}{i}\prod_{t=0}^{i-1}\frac{\alpha+c+t}{n+c+t}\Bigr)^{-1}\to1,
\]
the sum being finite. Since the posteriors of the compatible constituents
sum to \(1\), every \(C^w\) with \(w>c\) has posterior tending to \(0\). This is
(13); the formula depends on the sample only through \(n\) and \(c\), so it
applies along any sequence of samples in which the set of observed kinds
stays fixed.

For (14), let \(h\) have normal form (11), \(I_h\) indexing the constituents
that entail \(h\). Constituents are mutually exclusive, so
\(P(h/e)=\sum_{i\in I_h}P(C_i/e)\). If \(C^c\vdash h\), then \(C^c\) is a
disjunct and \(P(h/e)\geq P(C^c/e)\to1\). If not, every disjunct is a
constituent other than \(C^c\), and \(P(h/e)\leq1-P(C^c/e)\to0\), so
\(P(h/e)\not\to1\).

(3) A universal generalization declaring \(b\) given kinds empty is entailed
by exactly the constituents whose \(CT\) avoids those kinds. If \(h\) is
compatible with \(e\), none of the \(b\) kinds is observed, so the compatible
constituents entailing \(h\) are those adding \(i\) of the remaining
\(K-b-c\) kinds to the observed ones, \(\binom{K-b-c}{i}\) of width \(c+i\).
Summing (20) over them gives (23).

(4) If \(Q_i\) is observed (\(n_i\geq1\)), it belongs to every compatible
constituent, so (15\('\)) applies to each, and by total probability
\[
 P(Q_i(a_{n+1})/e)=\sum_{i'=0}^{K-c}\binom{K-c}{i'}P(C^{c+i'}/e)\frac{n_i+1}{n+c+i'} .
\]
Inserting (20) and using \(\Gamma(n+c+i')(n+c+i')=\Gamma(n+c+i'+1)\) gives
\[
 (n_i+1)\frac{\sum_{i'}\binom{K-c}{i'}\Gamma(\alpha+c+i')/\Gamma(n+c+i'+1)}
 {\sum_{i'}\binom{K-c}{i'}\Gamma(\alpha+c+i')/\Gamma(n+c+i')},
\]
which is (22) with \((x-1)!=\Gamma(x)\).
\end{proof}
\end{annexitem}

\begin{annexitem}{annex:p15-w02b-p15-030-priors}{thm:p15-w02b-p15-030-priors}{The prior assignments (17), (18), (19) and (25)}
\begin{proof}
(1) For \(w\geq1\) we have \(w\lambda/K>0\), and with \(\alpha=0\) the
weight is
\[
 \Gamma(w\lambda/K)/\Gamma(w\lambda/K)=1 .
\]
Equal weights give equal normalized priors.

(2) With \(G(w)=\Gamma(\alpha+w)/\Gamma(w)\),
\(G(w+1)/G(w)=(\alpha+w)/w>1\) for \(\alpha>0\). So \(G\), and hence
\(P(C^w)=G(w)/Z\), is strictly increasing in \(w\), which gives
\(P(C^w)<P(C^{w'})\iff w<w'\).

(3) For \(1\leq w<K\),
\[
 \frac{P(C^w)}{P(C^K)}=\frac{\Gamma(K)}{\Gamma(w)}\cdot\frac{\Gamma(\alpha+w)}{\Gamma(\alpha+K)}
 =\frac{\Gamma(K)}{\Gamma(w)}\prod_{t=w}^{K-1}\frac1{\alpha+t}\longrightarrow0
 \quad(\alpha\to\infty).
\]
There are finitely many consistent constituents and their priors sum to
\(1\), so \(P(C^K)=\bigl(1+\sum_{w<K}\binom Kw P(C^w)/P(C^K)\bigr)^{-1}\to1\),
and then each \(P(C^w)\leq1-P(C^K)\to0\) for \(w<K\).
\end{proof}
\end{annexitem}

\begin{annexitem}{annex:p15-w02b-p15-030-kemeny-carnap}{thm:p15-w02b-p15-030-kemeny-carnap}{The Kemeny--Carnap theorem on singular sentences}
\begin{proof}
Throughout, \(e\) is a sample description of the individuals
\(a_1,\ldots,a_n\), with counts \(n_1,\ldots,n_K\), and
\[
 f(n_i,n)=P(Q_i(a_{n+1})/e),
\]
by A5 the same \(f\) for every \(i\). Since the \(Q\)-predicates are exclusive and exhaustive, A1 gives
\begin{equation*}
 \textstyle\sum_{i=1}^Kf(n_i,n)=1\qquad(\ast)
\end{equation*}
for every count vector with sum \(n\). For \(n=0\), A4 makes all
\(f(0,0)\) equal, so \(f(0,0)=1/K\).

\emph{Step 1: \(f(\cdot,n)\) is affine for \(n\geq1\).} If \(K=2\) this is
assumed. Let \(K\geq3\) and \(0\leq x\leq n-1\). Compare the count vectors
\((x,1,n-x-1,0,\ldots,0)\) and \((x+1,0,n-x-1,0,\ldots,0)\): by \((\ast)\)
their \(f\)-sums are both \(1\), and they differ only in the first two
entries, so \(f(x,n)+f(1,n)=f(x+1,n)+f(0,n)\). Thus
\(f(x+1,n)-f(x,n)=f(1,n)-f(0,n)\) for \(0\leq x\leq n-1\), and
\(f(x,n)=a_n+b_nx\) with \(a_n=f(0,n)\), \(b_n=f(1,n)-f(0,n)\). Applying
\((\ast)\) to \((n,0,\ldots,0)\) gives
\begin{equation*}
 Ka_n+nb_n=1\qquad(\ast\ast).
\end{equation*}
In the case \(K=2\) with \(f(x,n)=a_n+b_nx\), \((\ast)\) for \((n,0)\) gives the
same \((\ast\ast)\).

\emph{Step 2: positivity.} Let \(n\geq1\) and let \(e\) put all \(n\)
individuals in \(Q_1\). If \(a_n=f(0,n)=0\), then
\(P(\neg Q_2(a_{n+1})/e)=1\), although \(e\) does not entail
\(\neg Q_2(a_{n+1})\); this contradicts A2. So \(a_n>0\) for all \(n\geq1\).

\emph{Step 3: exchangeability links \(n\) and \(n+1\).} With the same \(e\)
and \(j\neq1\), the permutation exchanging \(a_{n+1}\) and \(a_{n+2}\) fixes
\(e\), so by A3
\[
 P(Q_1(a_{n+1})\&Q_j(a_{n+2})/e)=P(Q_j(a_{n+1})\&Q_1(a_{n+2})/e).
\]
By the
product rule and A5 the left side is \(f(n,n)f(0,n+1)\) and the right side
\(f(0,n)f(n,n+1)\). Inserting the affine forms,
\((a_n+nb_n)a_{n+1}=a_n(a_{n+1}+nb_{n+1})\), that is
\(nb_na_{n+1}=na_nb_{n+1}\), so \(b_na_{n+1}=a_nb_{n+1}\) for \(n\geq1\). By
\((\ast\ast)\), \(a_n=(1-nb_n)/K\) and \(a_{n+1}=(1-(n+1)b_{n+1})/K\);
substituting and simplifying,
\begin{equation*}
 b_{n+1}(1+b_n)=b_n\qquad(n\geq1).\qquad(\dagger)
\end{equation*}

\emph{Step 4: solving.} If \(b_1=0\), then \((\dagger)\) gives \(b_2=0\), and
inductively \(b_n=0\) for all \(n\geq1\); then \(a_n=1/K\), so \(f\equiv1/K\)
(together with \(f(0,0)=1/K\)), and \(Kf(0,1)=1\).

If \(b_1\neq0\), put \(\lambda=1/b_1-1\). We show by induction that
\(b_n\neq0\) and \(1/b_n=n+\lambda\) for all \(n\geq1\). This holds for \(n=1\).
If it holds for \(n\), then \(b_n\neq-1\), since otherwise \((\dagger)\) would
read \(0=-1\); so \(b_{n+1}=b_n/(1+b_n)\neq0\) and
\(1/b_{n+1}=1/b_n+1=(n+1)+\lambda\). Then
\(a_n=(1-n/(n+\lambda))/K=\lambda/(K(n+\lambda))\). Choosing \(n\) with
\(n+\lambda>0\), Step 2 gives \(\lambda>0\). Hence \(n+\lambda>0\) for all \(n\),
and for \(0\leq x\leq n\),
\[
 f(x,n)=\frac{\lambda/K+x}{n+\lambda},
\]
which also gives \(f(0,0)=1/K\). Finally \(b_1=1-Ka_1=1-Kf(0,1)\), so
\(Kf(0,1)<1\) (as \(\lambda>0\) means \(b_1\in(0,1)\)) and
\(\lambda=1/(1-Kf(0,1))-1=Kf(0,1)/(1-Kf(0,1))\).

\emph{Converse.} For \(0<\lambda<\infty\) assign to each state description of
\(a_1,\ldots,a_N\) with counts \(N_1,\ldots,N_K\) the mass
\(\prod_j(\lambda/K)^{\overline{N_j}}/\lambda^{\overline N}\), where
\(x^{\overline m}=x(x+1)\cdots(x+m-1)\). This is the family of
\cref{thm:p15-carnap-history-family} with all \(\alpha_i=\lambda/K\): the masses
are positive, sum to \(1\), are exchangeable, agree under deleting the last
individual, and give predictive probability \((n_i+\lambda/K)/(n+\lambda)\).
Every singular sentence mentioning only \(a_1,\ldots,a_N\) is equivalent to a
disjunction of their state descriptions; let its probability be the sum of
their masses, which by the deletion property does not depend on the choice
of \(N\). This is a finitely additive probability on singular sentences
(A1). Every state description has positive mass, so \(P(h/e)=1\) forces every
state description consistent with \(e\) to entail \(h\), that is
\(\vdash e\supset h\) (A2). A permutation of individual constants carries
state descriptions to state descriptions with the same counts, hence the
same mass (A3); the masses are symmetric in the \(Q\)-predicates (A4); and
the predictive rule depends only on \((n_i,n)\) (A5). For \(\lambda=\infty\)
take independent uniform kinds, mass \(K^{-N}\) for each state description;
the same checks apply with predictive value \(1/K\). Finally \(\lambda=0\)
would give \(f(0,n)=0\) for \(n\geq1\), which Step 2 excludes.
\end{proof}
\end{annexitem}

\begin{annexitem}{annex:p15-w02b-p15-030-kdim}{thm:p15-w02b-p15-030-kdim}{Consequences of the stated \(K\)-dimensional formulas}
\begin{proof}
(1) In Carnap's continuum \(f(n_i,n)=(n_i+\lambda/K)/(n+\lambda)\) whatever
\(c\) is, so \(f(0,c,c)=(\lambda/K)/(c+\lambda)\). And
\((\lambda/K)/(c+\lambda)<1/K\iff\lambda<c+\lambda\), true as \(c\geq1\).

(2) By (26) each ratio \(\gamma_i/\delta_i\) lies in \((0,1]\), so their
product \(P(C^K)\) lies in \((0,1]\) and equals \(1\) iff every ratio equals
\(1\). If \(P(C^K)<1\), the consistent constituents, being exclusive and
jointly exhaustive, carry total probability \(1\), so those of width less
than \(K\) carry \(1-P(C^K)>0\).

(3) Evidence \(e\) in which all \(K\) kinds occur is compatible only with
\(C^K\), so \(P(C^K/e)=1\) and \(P(Q_i(a_{n+1})/e)=P(Q_i(a_{n+1})/e\&C^K)\). By
(27) with \(w=K\) this is \((n_i+\lambda/K)/(n+\lambda)\). For \(n=K+1\),
\(c=K\), \(n_i=1\) (one kind twice, the others once),
\(f(1,K+1,K)=(1+\lambda/K)/(K+1+\lambda)\). Then
\(Kf(1,K+1,K)=(K+\lambda)/(K+1+\lambda)\) and
\(1-Kf(1,K+1,K)=1/(K+1+\lambda)\), whose quotient is \(K+\lambda\);
subtracting \(K\) gives \(\lambda\).
\end{proof}
\end{annexitem}

\begin{annexitem}{annex:p15-w02b-p15-030-intersection}{thm:p15-w02b-p15-030-intersection}{When the two-dimensional likelihoods coincide with (27)}
\begin{proof}
Write \(x=\lambda_H/w\) and \(y=\lambda/K\), both positive; the identity reads
\((n_i+x)/(n+wx)=(n_i+y)/(n+wy)\). If \(x=y\) it holds for all \(n_i,n\).
Conversely, take \(n=1\), \(n_i=0\): \(x/(1+wx)=y/(1+wy)\), so
\(x(1+wy)=y(1+wx)\), that is \(x=y\), i.e. \(\lambda_H=w\lambda/K\). Applying
this for each \(w\) gives \(\lambda_H(w)=aw\) with \(a=\lambda/K\), the same for
all \(w\); with \(a=1\), \(\lambda_H(w)=w\) and the likelihood is
\((n_i+1)/(n+w)\), which is (15\('\)).
\end{proof}
\end{annexitem}

\begin{annexitem}{annex:p15-w02b-p15-030-uchii}{thm:p15-w02b-p15-030-uchii}{Uchii's convergence condition (31) and the Plenitude postulate}
\begin{proof}
Throughout, \(P\) gives S5-equivalent sentences equal probability, so if
\(A\vdash A'\) then \(P(A)=P(A\&A')\leq P(A')\). Two facts about
\(Q\)-predicates are used: they are pairwise exclusive and jointly
exhaustive, so for \(i\notin CT\), \(\vdash(x)\bigl[\bigvee_{j\in CT}Q_j(x)\bigr]\to(x)\sim Q_i(x)\)
and \(\vdash(x)\bigl[\bigwedge_{j\notin CT}\sim Q_j(x)\bigr]\to(x)\bigl[\bigvee_{j\in CT}Q_j(x)\bigr]\).
Write \(e\) for \(e^c_n\).

(1) Since \(P(C^c\&e)>0\), both \(P(C^c)>0\) and \(P(e)>0\). If
\(P(C^c\&B^c)=0\), then \(P(B^c\&C^c\&e)=0\) and \(P(B^c/C^c)=0\), so the
product identity holds with both sides \(0\). Otherwise the hypothesis gives
\[
 P(B^c\&C^c\&e)=P(e/C^c\&B^c)\,P(C^c\&B^c)=P(e/C^c)\,P(C^c\&B^c)
 =P(C^c\&e)\,\frac{P(C^c\&B^c)}{P(C^c)}=P(C^c\&e)\,P(B^c/C^c),
\]
and dividing by \(P(e)\) gives \(P(B^c\&C^c/e)=P(B^c/C^c)P(C^c/e)\). Now
\[
 P(B^c/e)=P(B^c/C^c)\,P(C^c/e)+P(B^c\&\sim C^c/e),\qquad
 0\leq P(B^c\&\sim C^c/e)\leq1-P(C^c/e).
\]
By (13), \(P(C^c/e^c_n)\to1\); the first term tends to \(P(B^c/C^c)\) and the
second to \(0\). Hence \(P(B^c/e^c_n)\to P(B^c/C^c)\), and the limit is \(1\)
iff \(P(B^c/C^c)=1\), which is (31).

(2) \(P(C^c)>0\) by part~(1), and \(P(B^c)=P(C^c)\), so
\(P(B^c/C^c)=P(B^c\&C^c)/P(C^c)=P(C^c\&B^c)/P(B^c)=P(C^c/B^c)\).

(3) First, for \(i\notin CT\), necessitation of the first fact and the K
axiom give \(B_{CT}\vdash\Box(x)\sim Q_i(x)\), which is S5-equivalent to
\(\sim\Diamond(\exists x)Q_i(x)\). If \(CT\neq CT'\), choose \(i\) in one set and
not the other, say \(i\in CT\setminus CT'\); then \(B_{CT}\vdash\Diamond(\exists x)Q_i(x)\)
and \(B_{CT'}\vdash\sim\Diamond(\exists x)Q_i(x)\), so \(B_{CT}\&B_{CT'}\) is
inconsistent. If \(i\in CT_0\setminus CT\), then \(C_{CT_0}\vdash(\exists x)Q_i(x)\)
while \(B_{CT}\vdash\Box(x)\sim Q_i(x)\vdash(x)\sim Q_i(x)\) by the T axiom of S5,
so \(B_{CT}\&C_{CT_0}\) is inconsistent. For satisfiability let
\(\emptyset\neq CT_0\subseteq CT\) (the empty constituent is itself
inconsistent), fix \(j_0\in CT_0\), and take the S5 model with worlds
\(w_0,w_1\), universal accessibility, and common domain
\(\{d_j\mid j\in CT\}\). Each \(Q\)-predicate fixes the sign of every primitive
predicate, so the primitive extensions at a world can be chosen to give
each individual any prescribed \(Q\)-predicate. At \(w_1\) let \(d_j\) satisfy
\(Q_j\); at \(w_0\) let \(d_j\) satisfy \(Q_j\) for \(j\in CT_0\) and \(Q_{j_0}\) for
\(j\in CT\setminus CT_0\). Then exactly the kinds in \(CT_0\) are exemplified
at \(w_0\), so \(C_{CT_0}\) is true there; at both worlds every individual
falls under some \(Q_j\) with \(j\in CT\), so \(\Box(x)\bigl[\bigvee_{j\in CT}Q_j(x)\bigr]\)
is true at \(w_0\); and each \(Q_j\), \(j\in CT\), is exemplified at \(w_1\), so
\(\Diamond(\exists x)Q_j(x)\) is true at \(w_0\). Thus \(B_{CT}\) is true at
\(w_0\) as well. The nomic constituents compatible with \(C^w=C_{CT_w}\) are
therefore those with \(CT\supseteq CT_w\); such \(CT\) is \(CT_w\) together with
an arbitrary subset of the \(K-w\) remaining indices, giving \(2^{K-w}\).

(4) For \(CT\neq CT_c\), part~(3) gives \(B_{CT}\vdash\sim B^c\), so
\(P(B_{CT}\&C^c)\leq P(\sim B^c\&C^c)=P(C^c)\bigl(1-P(B^c/C^c)\bigr)=0\), and
\(P(B_{CT}/C^c)=0\).

(5) Let \(i\notin CT_c\). By (10) and the first fact, \(C^c\vdash\sim(\exists x)Q_i(x)\).
If \(P(\Diamond(\exists x)Q_i(x))>0\), the postulate gives
\(P(\Diamond(\exists x)Q_i(x)\&\sim(\exists x)Q_i(x))=0\); if
\(P(\Diamond(\exists x)Q_i(x))=0\) the same equation holds trivially. Hence
\(P(C^c\&\Diamond(\exists x)Q_i(x))=0\). Put
\(N=\bigwedge_{i\notin CT_c}\sim\Diamond(\exists x)Q_i(x)\) (a tautology if \(c=K\)).
Then \(P(C^c\&\sim N)\leq\sum_{i\notin CT_c}P(C^c\&\Diamond(\exists x)Q_i(x))=0\),
so \(P(C^c\&N)=P(C^c)\). Next, \(C^c\&N\vdash B^c\): for \(i\in CT_c\),
\(C^c\vdash(\exists x)Q_i(x)\vdash\Diamond(\exists x)Q_i(x)\) by the T axiom;
and \(N\) is S5-equivalent to \(\bigwedge_{i\notin CT_c}\Box(x)\sim Q_i(x)\), hence
by K to \(\Box(x)\bigl[\bigwedge_{i\notin CT_c}\sim Q_i(x)\bigr]\), which by
necessitation of the second fact and K entails
\(\Box(x)\bigl[\bigvee_{i\in CT_c}Q_i(x)\bigr]\). Therefore
\(P(B^c\&C^c)\geq P(C^c\&N)=P(C^c)\), and \(P(B^c/C^c)=1\).
\end{proof}
\end{annexitem}

\begin{annexitem}{annex:p15-w02b-p15-030-information}{thm:p15-w02b-p15-030-information}{Identities among the information measures}
\begin{proof}
All logarithms are of positive numbers: \(P(h\&e)>0\) gives \(P(h),P(e)>0\).

(1) We have
\[
 \operatorname{inf}(h\&e)-\operatorname{inf}(e)=-\log P(h\&e)+\log P(e)
 =-\log\frac{P(h\&e)}{P(e)}=-\log P(h/e).
\]

(2) \(\operatorname{transinf}(h/e)=-\log P(h)+\log P(h/e)\).
\(\operatorname{transcont}_{add}(h/e)=(1-P(h))-[(1-P(h\&e))-(1-P(e))]
=1-[P(h)+P(e)-P(h\&e)]=1-P(h\lor e)\).
\(\operatorname{transcont}_{cond}(h/e)=(1-P(h))-(1-P(h/e))=P(h/e)-P(h)\).

(3) By (2) with \(h\) and \(e\) exchanged,
\(\operatorname{transinf}(e/h)=\log P(e/h)-\log P(e)\). Dividing by
\(\operatorname{inf}(e)=-\log P(e)\), which is nonzero because \(P(e)<1\), gives
\[
 \frac{\log P(e)-\log P(e/h)}{\log P(e)} .
\]
Next,
\(\operatorname{transcont}_{add}(e/h)=1-P(h\lor e)=P(\sim h\&\sim e)\), and
dividing by \(\operatorname{cont}(e)=P(\sim e)>0\) gives \(P(\sim h/\sim e)\).
Finally \(\operatorname{transcont}_{cond}(e/h)=P(e/h)-P(e)\), divided by \(1-P(e)\).

(4) \(\log\) is strictly increasing, so \(\log P(h/e)-\log P(h)>0\iff
P(h/e)>P(h)\iff P(h/e)-P(h)>0\). And \(1-P(h\lor e)>0\iff P(h\lor e)<1\).

(5) If \(h\vdash e\), then \(P(e/h)=1\) and \(P(h\lor e)=P(e)\), so the three
quotients in (3) are \(\log P(e)/\log P(e)\), \((1-P(e))/(1-P(e))\) and
\((1-P(e))/(1-P(e))\), each equal to \(1\).
\end{proof}
\end{annexitem}

\begin{annexitem}{annex:p15-w02b-p15-030-relevance}{thm:p15-w02b-p15-030-relevance}{Positive relevance, the principal monotonicities and their hypotheses}
\begin{proof}
(1) Multiplying by \(P(e)>0\): \(P(h/e)>P(h)\iff P(h\&e)>P(h)P(e)\).
Dividing instead by \(P(h)>0\): \(\iff P(e/h)>P(e)\). For the third form,
\(P(h)=P(e)P(h/e)+P(\sim e)P(h/\sim e)\) with \(P(e)+P(\sim e)=1\) and both
positive, so \(P(h)\) lies strictly between \(P(h/e)\) and \(P(h/\sim e)\)
unless they are equal; hence \(P(h/e)>P(h)\iff P(h/e)>P(h/\sim e)\).

(2) If \(h\vdash e\), \(P(h\&e)=P(h)\) and \(P(h/e)=P(h)/P(e)>P(h)\), since
\(P(h)>0\) and \(P(e)<1\).

(3) If \(e\vdash h_1\supset h_2\), then \(e\&h_1\&\sim h_2\) is inconsistent, so
\(P(h_1/e)=P(h_1\&h_2/e)+P(h_1\&\sim h_2/e)=P(h_1\&h_2/e)\leq P(h_2/e)\).

(4) Let \(C\) be a disjunct of \(h_1\) consistent with \(e\). Then
\(C\vdash h_1\), so \(C\&e\vdash h_2\). A constituent decides every
generalization of \(L\), so \(C\vdash h_2\) or \(C\vdash\sim h_2\); the latter
would make \(C\&e\) inconsistent. So \(C\vdash h_2\), i.e. \(C\) is a disjunct
of \(h_2\) consistent with \(e\). The minimum over a larger set is no larger,
so \(\operatorname{corr}^e(h_1,e)\geq\operatorname{corr}^e(h_2,e)\). If
\(\vdash h_1\supset h_2\), every disjunct of \(h_1\) entails \(h_2\), so
\(I_{h_1}\subseteq I_{h_2}\) and the same comparison holds for the minimum over
all disjuncts.

(5) Take \(K=2\), \(e=Q_1(a_1)\), and let \(C_1\) say that only \(Q_2\) is
exemplified and \(C_2\) that only \(Q_1\) is; \(C_1\) is inconsistent with
\(e\), so \(P(C_1/e)=0\). Let \(P\) be any measure with \(P(C_2/e)>0\) (for
example any system of \cref{thm:p15-w02b-p15-030-combined}). Put
\(h_1=C_1\lor C_2\) and \(h_2=C_2\). Since \(e\vdash\sim C_1\), \(e\vdash
h_1\supset h_2\). But \(\operatorname{corr}(h_1,e)=\min(0,P(C_2/e))=0<
P(C_2/e)=\operatorname{corr}(h_2,e)\).
\end{proof}
\end{annexitem}

\begin{annexitem}{annex:p15-w02b-p15-030-utilities}{thm:p15-w02b-p15-030-utilities}{Expected epistemic utilities (48), (50), (52), (53), (55)}
\begin{proof}
Write \(x=P(g/e)=\sum_{j\in I_g}P(h_j/e)\). For a utility taking value
\(u_1\) when \(g\) is true and \(u_0\) when false, (46) gives
\(U(g/e)=xu_1+(1-x)u_0\).

(1) \(u_1=1\), \(u_0=0\) gives \(x\) (48). If \(u_1=u_0=1-P(g)\), then
\(U=1-P(g)\).

(2) \(u_1=1+qI(g)\), \(u_0=qI(g)\) gives \(x+qI(g)\) (50).

(3) \(u_1=1-P(g)\), \(u_0=-(1-P(\sim g))=-P(g)\):
\(U=x(1-P(g))-(1-x)P(g)=x-P(g)\) (52).

(4) Put \(y=P(g/\sim e)\). \(u_1=1-y\), \(u_0=-y\):
\(U=x(1-y)-(1-x)y=x-y\). Since \(P(g)=P(e)x+P(\sim e)y\),
\(x-P(g)=(1-P(e))x-P(\sim e)y=P(\sim e)(x-y)\), so
\(U=[x-P(g)]/P(\sim e)\).

(5) \(u_1=1-P(g)\), \(u_0=-qP(g)\):
\(U=x(1-P(g))-(1-x)qP(g)=x-qP(g)-(1-q)P(g)x\). The last term vanishes iff
\(q=1\) or \(P(g)x=0\). With \(u_1=1-qP(g)\) instead,
\(U=x(1-qP(g))-(1-x)qP(g)=x-qP(g)\).
\end{proof}
\end{annexitem}

\begin{annexitem}{annex:p15-w02b-p15-030-analogy}{thm:p15-w02b-p15-030-analogy}{The analogy formulas (58)--(60)}
\begin{proof}
\(F_1\) is the disjunction of the \(w_1=2^{k-m}\) \(Q\)-predicates whose first
\(m\) signs are positive, and \(F_2\) of the \(w_2=2^{k-n}\) whose first \(n\)
signs are positive; the latter set is contained in the former, so
\(F_2\vdash F_1\).

By A1 and A4, \(P(Q_i(b))=1/K\) for each \(i\), so \(P(F_r(b))=w_r/K\) and
\(P(F_2(b)/F_1(b))=P(F_2(b))/P(F_1(b))=w_2/w_1\).

By A3 we may take \(a=a_1\) and \(b=a_2\). By A6 applied to the evidence
\(Q_i(a)\) (\(n=1\), \(c=1\)),
\(P(Q_j(b)/Q_i(a))\) is \(f(1,1,1)\) for \(j=i\) and \(\varphi=f(0,1,1)\) for
\(j\neq i\); by A1, \(f(1,1,1)+(K-1)\varphi=1\). For \(i\) in \(F_2\) (hence in
\(F_1\)),
\[
 P(F_r(b)/Q_i(a))=f(1,1,1)+(w_r-1)\varphi=1-(K-w_r)\varphi\qquad(r=1,2).
\]
Summing over the \(w_2\) kinds of \(F_2\) with weights \(P(Q_i(a))=1/K\),
\[
 P(F_r(b)\&F_2(a))=\frac{w_2}{K}\bigl(1-(K-w_r)\varphi\bigr).
\] By A2,
\(f(1,1,1)>0\) (otherwise \(P(\neg Q_i(b)/Q_i(a))=1\) without entailment), so
\(1-(K-w_1)\varphi\geq f(1,1,1)>0\). As \(F_2(b)\) entails \(F_1(b)\),
\[
 P(F_2(b)/F_1(b)\&F_2(a))=\frac{P(F_2(b)\&F_2(a))}{P(F_1(b)\&F_2(a))}
 =\frac{1-(K-w_2)\varphi}{1-(K-w_1)\varphi}.
\]
This exceeds \(w_2/w_1\) iff
\(w_1(1-(K-w_2)\varphi)>w_2(1-(K-w_1)\varphi)\), i.e.
\((w_1-w_2)-\varphi(Kw_1-Kw_2)>0\), i.e. \((w_1-w_2)(1-K\varphi)>0\), i.e.
\(K\varphi<1\) since \(w_1>w_2\). And \(f(1,1,1)=1-(K-1)\varphi>\varphi\iff
K\varphi<1\).

With \(\varphi=(\lambda/K)/(1+\lambda)\),
\(1-(K-w)\varphi=[(1+\lambda)-(K-w)\lambda/K]/(1+\lambda)=(1+w\lambda/K)/(1+\lambda)\),
so the quotient is \((1+w_2\lambda/K)/(1+w_1\lambda/K)\), (59); at \(\lambda=K\)
it is \((1+w_2)/(1+w_1)\), (60).
\end{proof}
\end{annexitem}

\begin{annexitem}{annex:p15-w02b-p15-030-truthlikeness}{thm:p15-w02b-p15-030-truthlikeness}{Properties of \(Tr\) and \(\operatorname{ver}\)}
\begin{proof}
Write \(*\) for the index of \(h^*\) and \(S=\sum_{j\in I}\Delta_{*j}\), which is
positive because \(m\geq2\) and \(\Delta_{*j}>0\) for \(j\neq*\).

(1) \(Tr(g,h^*)=1\) iff \(\gamma\Delta_{\min}(h^*,g)+\gamma'\Delta_{\mathrm{sum}}(h^*,g)=0\)
iff both terms vanish (they are nonnegative and \(\gamma,\gamma'>0\)).
\(\Delta_{\min}(h^*,g)=0\) iff \(*\in I_g\); \(\Delta_{\mathrm{sum}}(h^*,g)=0\) iff
\(\Delta_{*j}=0\) for all \(j\in I_g\), iff \(I_g\subseteq\{*\}\). Together:
\(I_g=\{*\}\), i.e. \(\vdash g\equiv h^*\).

(2) For a tautology \(I_t=I\): \(\Delta_{\min}=\Delta_{**}=0\) and
\(\Delta_{\mathrm{sum}}=S/S=1\), so \(Tr(t,h^*)=1-\gamma'\). For \(g=h_i\):
\(\Delta_{\min}=\Delta_{*i}\), \(\Delta_{\mathrm{sum}}=\Delta_{*i}/S\), so
\(Tr(h_i,h^*)=1-(\gamma+\gamma'/S)\Delta_{*i}\), strictly decreasing in
\(\Delta_{*i}\).

(3) With \(\Delta_{ij}=1\) for \(i\neq j\): \(S=m-1\),
\(\Delta_{\min}(h^*,g)=1-\operatorname{tv}(g)\) and
\(\Delta_{\mathrm{sum}}(h^*,g)=(|I_g|-\operatorname{tv}(g))/(m-1)\). Substituting
\(|I_g|=m-(m-1)I(g)\),
\[
 Tr(g,h^*)=\Bigl[1-\gamma-\frac{\gamma'm}{m-1}\Bigr]
 +\Bigl(\gamma+\frac{\gamma'}{m-1}\Bigr)\operatorname{tv}(g)+\gamma'I(g)
 =\kappa+\mu\bigl[\operatorname{tv}(g)+qI(g)\bigr]
\]
with \(\mu=\gamma+\gamma'/(m-1)>0\) and \(q=\gamma'/\mu>0\). Then
\(q\leq1\iff\gamma'\leq\gamma+\gamma'/(m-1)\iff\gamma'(m-2)\leq\gamma(m-1)\).

(4) \(\operatorname{tv}(g,h_i)\) is \(1\) for \(i\in I_g\) and \(0\) otherwise, so
\(\sum_iP(h_i/e)\operatorname{tv}(g,h_i)=\sum_{i\in I_g}P(h_i/e)=P(g/e)\). If
\(P(h_j/e)=1\), all other \(P(h_i/e)\) are \(0\) and (70) reduces to
\(Tr(g,h_j)\).

(5) \(B\) is finite and each \(Tr(g,h_i)\in[1-\gamma-\gamma',1]\) is fixed. By
(13), \(P(C^c/e)\to1\) and every other \(P(h_i/e)\to0\), so the finite sum (70)
tends to \(Tr(g,C^c)\), which is (72). Hence \(\operatorname{ver}(g/e)\to1\) iff
\(Tr(g,C^c)=1\), iff \(\vdash g\equiv C^c\) by part~(1) with \(C^c\) in the role
of the target, which is (73).
\end{proof}
\end{annexitem}

\begin{annexitem}{annex:p15-w02b-p15-030-pat}{thm:p15-w02b-p15-030-pat}{Probable approximate truth and convergence under ESC}
\begin{proof}
(1) For \(i\in I_g\), \(\Delta_{\min}(h_i,g)\leq\Delta_{ii}=0\), so
\(I_g\subseteq V_\varepsilon(g)\) for every \(\varepsilon\geq0\), and
\(P(g/e)\leq\sum_{h_i\in V_\varepsilon(g)}P(h_i/e)\). If \(\varepsilon\leq\varepsilon'\),
then \(V_\varepsilon(g)\subseteq V_{\varepsilon'}(g)\), so \(PAT\) is nondecreasing in
\(\varepsilon\). For \(\varepsilon=0\), \(\Delta_{\min}(h_i,g)\leq0\) iff \(\Delta_{ij}=0\)
for some \(j\in I_g\) iff \(i\in I_g\); so \(V_0(g)=I_g\) and \(PAT_1(g/e)=P(g/e)\).

(2) By (75) and the definition of \(g^\varepsilon\),
\(PAT_{1-\varepsilon}(g/e)=P(g^\varepsilon/e)\). For an instance with
\(P(g/e)=0\): \(B=\{h_1,h_2\}\), \(\Delta_{12}=0.1\), \(g=h_1\), \(\varepsilon=0.1\),
\(P(h_2/e)=1\); then \(V_\varepsilon(g)=\{h_1,h_2\}\) and \(PAT_{0.9}(g/e)=1\).
For probable verisimilitude, by \cref{def:p15-w02b-p15-030-nn-pv} its value is
the posterior mass of \(\{h_i\mid Tr(g,h_i)\geq1-\varepsilon\}\). In the
instance of \cref{def:p15-w02b-p15-030-nn-pat-worked}, the values
\(Tr(h_1,h_1)=1\), \(Tr(h_1,h_2)=47/55\) and \(Tr(h_1,h_3)=-1/4\) follow from
part~(2) of \cref{thm:p15-w02b-p15-030-truthlikeness} (with sums of distances
\(1.1\) and \(2\) for the targets \(h_2\) and \(h_3\)); so the set is
\(\{h_1,h_2\}\), and \(P(Tr(g,h^*)\geq0.8\,/\,e)=0+0.7>0\) while \(P(g/e)=P(h_1/e)=0\).

(3) If \(g_1\vdash g_2\), then \(I_{g_1}\subseteq I_{g_2}\), so
\(\Delta_{\min}(h_i,g_2)\leq\Delta_{\min}(h_i,g_1)\) for every \(i\), hence
\(V_\varepsilon(g_1)\subseteq V_\varepsilon(g_2)\) and the sums compare.

(4) Under ESC the evidence is true, so every kind it exhibits exists, and
every existing kind is exhibited; the set of observed kinds is therefore
exactly the set of nonempty kinds, and the constituent \(C^c\) claiming
exactly the observed kinds is the true constituent \(C^*\). Substituting
\(C^*\) for \(C^c\) in (13) and in parts~(5) of
\cref{thm:p15-w02b-p15-030-truthlikeness} gives (13\('\)), (72\('\)) and
(73\('\)), under the same hypotheses on the measure that yield (13).

(5) \(B\) is finite. By (13), \(P(C^c/e)\to1\); since the posteriors of the
other members of \(B\) are nonnegative and sum to \(1-P(C^c/e)\), each tends to
\(0\). A sum of posteriors over a fixed subset \(A\subseteq B\) therefore tends to
\(1\) if \(C^c\in A\) and to \(0\) if not. The sets \(V_\varepsilon(g)\) and
\(\{h_i\mid Tr(g,h_i)\geq1-\varepsilon\}\) depend only on \(g\), \(\varepsilon\),
\(\Delta\), \(\gamma\) and \(\gamma'\), not on \(n\); applying this to them, with
(75) and \cref{def:p15-w02b-p15-030-nn-pv}, gives both claims, the first
membership condition being \(C^c\in V_\varepsilon(g)\) and the second
\(Tr(g,C^c)\geq1-\varepsilon\).

(6) Under ESC, part~(4) gives \(C^c=C^*\). Substituting in part~(5):
\(PAT_{1-\varepsilon}(g/e)\to1\) iff \(C^*\in V_\varepsilon(g)\), iff
\(\Delta_{\min}(C^*,g)\leq\varepsilon\) by (74), which by (74) and the
definition of \(g^\varepsilon\) says that \(g^\varepsilon\) is true, i.e.\ that
\(g\) is approximately true within \(\varepsilon\); otherwise the limit is
\(0\). Likewise \(P(Tr(g,h^*)\geq1-\varepsilon\,/\,e)\to1\) iff
\(Tr(g,C^*)\geq1-\varepsilon\), and \(\to0\) otherwise.
\end{proof}
\end{annexitem}
